\chapter{Kiến trúc hệ thống (System Architecture)}

\section{Sơ đồ luồng dữ liệu (Data Flow)}
Quy trình xử lý dữ liệu của hệ thống trong giai đoạn hiện tại được thiết kế theo dạng đường ống (pipeline) tuyến tính, được minh họa trong Hình \ref{fig:data-flow}:

\begin{figure}[H]
\centering
\begin{tikzpicture}[
    node distance=0.6cm and 0.3cm,
    box/.style={rectangle, draw, rounded corners, minimum height=1cm, minimum width=2.8cm, text centered, font=\small},
    data/.style={box, fill=blue!15, text width=2.6cm},
    process/.style={box, fill=orange!20, text width=2.6cm},
    storage/.style={box, fill=green!15, text width=2.6cm},
    user/.style={box, fill=red!12, text width=2.6cm},
    arrow/.style={->, thick, >=stealth},
    label/.style={font=\scriptsize\itshape, text=gray}
]

% Row 1: Data pipeline
\node[data] (raw) {Raw Data\\{\scriptsize arXiv JSON 3GB+}};
\node[process, right=of raw] (clean) {Data Cleaning\\{\scriptsize Python streaming}};
\node[data, right=of clean] (json) {Clean JSON\\{\scriptsize 100k docs, 350MB}};
\node[process, right=of json] (bulk) {Bulk Ingestion\\{\scriptsize batch 5,000}};
\node[storage, right=of bulk] (es) {Elasticsearch\\{\scriptsize Inverted Index}};

% Row 2: Query pipeline
\node[user, below=1.2cm of raw] (user) {User Query\\{\scriptsize keyword + filter}};
\node[process, right=of user] (qproc) {Query Processing\\{\scriptsize tokenize, analyze}};
\node[process, right=of qproc] (bm25) {BM25 Retriever\\{\scriptsize scoring + ranking}};
\node[process, right=of bm25] (filter) {Filter\\{\scriptsize year, category}};
\node[user, right=of filter] (result) {Results\\{\scriptsize top-k ranked}};

% Arrows - data pipeline
\draw[arrow] (raw) -- (clean);
\draw[arrow] (clean) -- (json);
\draw[arrow] (json) -- (bulk);
\draw[arrow] (bulk) -- (es);

% Arrows - query pipeline
\draw[arrow] (user) -- (qproc);
\draw[arrow] (qproc) -- (bm25);
\draw[arrow] (bm25) -- (filter);
\draw[arrow] (filter) -- (result);

% Connection between two pipelines
\draw[arrow, dashed, blue!60] (es.south) -- ++(0,-0.35) -| (bm25.north);

% Labels
\node[label, above=0.15cm of clean] {data\_prep.py};
\node[label, above=0.15cm of bulk] {index\_data.py};
\node[label, below=0.15cm of qproc] {search.py};

\end{tikzpicture}
\caption{Sơ đồ luồng dữ liệu (Data Flow Pipeline) của hệ thống}
\label{fig:data-flow}
\end{figure}

\noindent Quy trình bao gồm hai luồng chính:

\textbf{Luồng nạp dữ liệu (Data Ingestion Pipeline):}
\begin{enumerate}
    \item \textbf{Raw Data:} Dữ liệu gốc định dạng JSON từ arXiv (3GB+, line-delimited).
    \item \textbf{Data Cleaning:} Script Python đọc từng dòng (streaming), lọc bỏ các bản ghi lỗi, và trích xuất tập con 100,000 bài báo thuộc lĩnh vực Computer Science.
    \item \textbf{Data Ingestion:} Dữ liệu sạch được đóng gói thành các batch 5,000 documents và gửi qua Elasticsearch Bulk API.
    \item \textbf{Indexing \& Storage:} Elasticsearch phân tích văn bản (tokenization) và tạo Inverted Index cho các trường \texttt{text}, đồng thời lưu trữ nguyên bản các trường \texttt{keyword}.
\end{enumerate}

\textbf{Luồng truy vấn (Query Pipeline):}
\begin{enumerate}
    \item \textbf{User Query:} Người dùng nhập câu truy vấn từ khóa, kèm theo bộ lọc (năm, danh mục) nếu cần.
    \item \textbf{Query Processing:} Hệ thống phân tích cú pháp (tokenize) câu truy vấn.
    \item \textbf{BM25 Retrieval:} Đối chiếu với Inverted Index, tính điểm BM25 cho từng document.
    \item \textbf{Filtering \& Ranking:} Áp dụng bộ lọc exact-match và trả về top-k kết quả được xếp hạng theo độ liên quan.
\end{enumerate}

\section{Kiến trúc hạ tầng (Infrastructure)}

\subsection{Công nghệ sử dụng}
\begin{itemize}
    \item \textbf{Python 3.x:} Xử lý dữ liệu (Data Processing) và tương tác với Elasticsearch thông qua thư viện \texttt{elasticsearch-py}.
    \item \textbf{Elasticsearch 8.12:} Search Engine cốt lõi, hỗ trợ Inverted Index, BM25 scoring, và filtering \cite{elasticsearch_docs}.
    \item \textbf{Docker \& Docker Compose:} Đóng gói và triển khai toàn bộ hạ tầng (Elasticsearch + Kibana) trong các container độc lập \cite{docker_docs}.
    \item \textbf{Kibana 8.12:} Giao diện quản trị và trực quan hóa dữ liệu, hỗ trợ Dev Tools để thử nghiệm truy vấn.
\end{itemize}

\subsection{Mô hình triển khai}
Hệ thống được triển khai dưới dạng Single-node cluster trên môi trường Docker cục bộ, sử dụng cấu hình Docker Compose như sau:

\begin{lstlisting}[caption={Docker Compose -- Elasticsearch \& Kibana},label={lst:docker-compose}]
services:
  elasticsearch:
    image: elasticsearch:8.12.2
    environment:
      - discovery.type=single-node
      - xpack.security.enabled=false
      - "ES_JAVA_OPTS=-Xms2g -Xmx2g"
    ports:
      - "9200:9200"
    volumes:
      - es_data:/usr/share/elasticsearch/data

  kibana:
    image: kibana:8.12.2
    environment:
      - ELASTICSEARCH_HOSTS=http://elasticsearch:9200
    ports:
      - "5601:5601"
\end{lstlisting}

\subsection{Đặc tính Big Data}
Mặc dù chạy trên single-node, kiến trúc bên trong của Elasticsearch vẫn duy trì các khái niệm phân tán như \textbf{Index}, \textbf{Shard} (phân mảnh dữ liệu) và \textbf{Replica} (bản sao). Điều này cho phép hệ thống dễ dàng mở rộng (scale-out) bằng cách thêm các node mới vào cluster khi khối lượng dữ liệu tăng lên mức hàng triệu bài báo mà không cần thay đổi kiến trúc ứng dụng.

\section{Index Mapping \& Bulk Strategy}

\subsection{Mapping Definition}
Lược đồ dữ liệu (Mapping) được định nghĩa tường minh để tối ưu hóa hiệu năng cho từng loại truy vấn:

\begin{lstlisting}[caption={Elasticsearch Index Mapping},label={lst:mapping}]
{
  "mappings": {
    "properties": {
      "id":         { "type": "keyword" },
      "title":      { "type": "text" },
      "abstract":   { "type": "text" },
      "categories": { "type": "keyword" },
      "authors":    { "type": "text" },
      "year":       { "type": "keyword" }
    }
  },
  "settings": {
    "number_of_shards": 1,
    "number_of_replicas": 0
  }
}
\end{lstlisting}

Trong đó:
\begin{itemize}
    \item Các trường \texttt{title}, \texttt{abstract}, \texttt{authors} được map dưới dạng \texttt{text} -- Elasticsearch tự động thực hiện phân tích cú pháp (tokenization, lowercasing, stop-word removal) và xây dựng Inverted Index để phục vụ tìm kiếm BM25.
    \item Các trường \texttt{id}, \texttt{year}, \texttt{categories} được map dưới dạng \texttt{keyword} -- lưu trữ nguyên bản, không qua phân tích, phục vụ exact-match filtering với tốc độ O(1).
\end{itemize}

\subsection{Bulk Indexing Strategy}
Thay vì gửi từng tài liệu qua HTTP POST riêng lẻ (rất chậm và tốn tài nguyên mạng), hệ thống sử dụng hàm \texttt{helpers.bulk()} của thư viện \texttt{elasticsearch-py}. Dữ liệu được gom thành các lô (batch) với kích thước \textbf{5,000 documents} mỗi lần gửi, giúp:
\begin{itemize}
    \item Giảm số lượng HTTP round-trip từ 100,000 xuống còn 20 requests.
    \item Tối ưu hóa băng thông mạng và giảm tải cho bộ nhớ heap của Elasticsearch.
    \item Toàn bộ 100,000 documents được index trong vòng vài chục giây.
\end{itemize}